The Mapper algorithm has become a popular tool in Topological Data Analysis 
(TDA) since its invention by Singh, et al. in \cite{Singh2007}. Mapper takes 
as input a dataset and outputs a simplicial complex. The raw input, 
interpreted as a point-cloud, lacks any meaningful topological structure. 
The resulting simplicial complex (mapper-complex) is an approximation 
for the topology of the underlying space from which the data may have been 
sampled.

Mapper is not unique in this regard. There are many simplicial-complex 
constructions in TDA which accomplish this (Vietoris-Rips complexes, 
Alpha complexes, etc.). However the main benefit (and curse) of mapper is 
the wide-variety of hyperparameters involved in its construction. A 
principled and informed application of Mapper can reveal interesting 
structures and analysis from one's data. However, it has been shown in 
\cite{alvarado2024graphmappergraph} that given any sufficiently large 
dataset and graph $G$ there is a choice of hyperparameters for which mapper's 
output is isomorphic to $G$. So an unprincipled approach can very well 
yield misleading conclusions.

There are as of yet no educational tools for the Mapper algorithm which 
are interactive and encourage a principled approach to the mapper algorithm.
ShinyMappeR is a software which aims to fill this gap for educators, students, 
and practitioners. ShinyMappeR is written in R as a shinyApp and can be used 
in the browser. ShinyMappeR gives visual feedback and provides example 
datasets which showcase some common issues that may arise in parameter selection 
for the mapper algorithm. ShinyMappeR uses the MappeR implementation of the 
mapper algorithm and is extensible to demonstrate new research methods. 
For example: G-Mapper comes equipped as a covering scheme option and allows for live 
feedback for how changing hyperparameters impacts every aspect of the mapper 
graph construction.

We now outline the rest of this paper. Section 2 continues with a detailed
description of the mapper algorithm and its motivations. Section 3 goes over common 
traps when implementing mapper and how shinymapper example datasets can teach 
these. Section 4 gives a breakdown of the features and architecture of shinyMappeR software. 
And section 5 outlines extensiblity for new research techniques.
